\documentclass{article}
\usepackage{graphicx} % Required for inserting images

\title{Upgraded SW Models}
\date{October 2024}

\begin{document}

\maketitle

\section{Shallow Water Equations with temperature gradient}
Integrating the primitive equations vertically, we get that for a layer going from $z_1$ to $z_2$, with vertically averaged density being $\bar{\rho}$:
\begin{equation}
    \bar{\rho}\left(z_2-z_1\right)\left(D_t\mathbf{v}+\mathbf{f}\times\mathbf{v}\right)=-\nabla\left(-g\bar{\rho}\frac{(z_2-z_1)^2}{2}+\left(z_2-z_1\right)p_{z_1}\right)-\nabla z_1 p_{z_1}+\nabla z_2 p_{z_2}
\end{equation}
Note that here we used
\begin{equation}
    p_z=-g\bar{\rho}\left(z-z_1\right)+p_{z_1}
\end{equation}
We can write it in terms of geopotential $\phi=-g\frac{\bar{\rho}}{\rho_0}(z-z_1)+\phi_{z_1}$ and buoyancy as
\begin{equation}
    \left(z_2-z_1\right)\left(D_t\mathbf{v}+\mathbf{f}\times\mathbf{v}\right)=-\nabla\left(-g\frac{\bar{\rho}}{\rho_0}\frac{(z_2-z_1)^2}{2}+\left(z_2-z_1\right)\phi_{z_1}\right)-\nabla z_1 \phi_{z_1}+\nabla z_2 \phi_{z_2}
\end{equation}
The integration of continuity equation, combined with incompressibility condition $\nabla \cdot \mathbf{v}=0$ and mean field hypothesis $\left<\rho \mathbf{v}\right>=\left<\rho \right>\left<\mathbf{v}\right>$, similar to normal SW, becomes
\begin{equation}
    \partial_t\left(z_2-z_1\right)+\nabla\cdot\left((z_2-z_1)\mathbf{v}\right)=0
\end{equation}
And the column-averaged density is a materially conserved tracer
\begin{equation}
    D_t \bar{\rho} = 0
\end{equation}

\section{Adding Moisture into the SW Equations}
Again we start from primitive equations:
\begin{equation}
    D_t \mathbf{v}+\mathbf{f}\times\mathbf{v}=-\nabla\Phi
\end{equation}
\begin{equation}
    \nabla \mathbf{v}+\partial_z w=0
\end{equation}
\begin{equation}
    D_t(\theta+\frac{L}{C_p}q)=0
\end{equation}
And if there is no condensation, $\theta$ and $q$ are separately conserved. Now considering a two-layer structure, with the bottom layer based at $z_0$ and top layer at $z_1$, top at $z_2$.
\begin{equation}
    w_0=\frac{dz_0}{dt}, w_1=\frac{dz_1}{dt}+W_1, w_2=\frac{dz_2}{dt}+W_2
\end{equation}
So the continuity equations write
\begin{equation}
    \partial_t h_1 + \nabla\cdot(h_1\mathbf{v}_1)=W_0-W_1
\end{equation}
\begin{equation}
    \partial_t h_2 + \nabla\cdot(h_2\mathbf{v}_2)=W_1-W_2
\end{equation}
The vertically averaged momentum equation then becomes
\begin{equation}
    \partial_t \mathbf{V}_1+(\mathbf{v}_1\cdot\nabla)\mathbf{V}_1+\mathbf{f}\times\mathbf{V}_1=-\nabla \Phi|_{z_1}+g\frac{\theta_1}{\theta_0}\nabla z_1
\end{equation}
\begin{equation}
    \partial_t \mathbf{V}_2+(\mathbf{V}_2\cdot\nabla)\mathbf{V}_2+\mathbf{f}\times\mathbf{V}_2=-\nabla \Phi|_{z_2}+g\frac{\theta_2}{\theta_0}\nabla z_2+\frac{v_1-v_2}{h}W_1
\end{equation}
Here if we take $\phi_0=0$, $\phi_1=g\frac{\theta_1}{\theta_0}(z_1-z_0)$, and $\phi_2=g\frac{\theta_1}{\theta_0}(z_1-z_0)+g\frac{\theta_2}{\theta_0}(z_2-z_1)$. Taking the vertically integrated water content in layer $i$ be $Q_i$, we have a local conservation
\begin{equation}
    \partial_t Q_i+\nabla\cdot(Q_i\mathbf{V}_i)=E_i-P_i
\end{equation}
Here $E$ and $P$ are the evaporation and precipitation respectively, which feeds in the forcing of the model. The vertical speed at the interfaces are directly related to the precipitation rate
\begin{equation}
    W_i=\beta_i P_i, \beta_i=\frac{L}{C_p(\theta_2-\theta_1)}
\end{equation}
Under the assumpation of a stably stratified background and a dry layer 2, only layer 1 has precipitations, we have
\begin{equation}
    \beta_1\approx \frac{1}{q_1}>0
\end{equation}
Here we can use a relaxation scheme to calculate the precipitation rate, using the saturated $Q_s=q_s(P,T)h$, or $P_1=\tau_1^{-1}(Q-Q_s)$, we use a realistic value of $P$ and $T$ to obtain the saturated vapor mixing ratio at this point. In addition, we need a water mixing ratio $q_0$, which is assumed to be a uniform distributed in the abyssal layer. What is the speed that precipitation is evaporated and fed back into the weather layer again, or, how do we calculate $W_0$? We assume all the precipitation falls into the abyssal layer and averaged, and that the evaporation rate is proportional to this precipitation,
\begin{equation}
    \frac{d\alpha_0}{dt}=\bar{P}_1 - \tau_2^{-1}\alpha_0
\end{equation}
And the speed means the same amount of the moisture is fed back to layer 1 from layer 0, so that
\begin{equation}
    W_0=\frac{\alpha_0}{\tau_2 q_0}
\end{equation}
This also gives the value of $E$, being
\begin{equation}
    E=W_0q_0
\end{equation}
Thus we close the system. Still, in addition to an infinitely deep lower layer (abyssal layer), we would like to note that the upper layer has very small impact on the lower layer, especially when we ignore the momentum transfer between these two layers. Ignoring the highest layer, just with the weather layer, we can have a one-layer model for the water. The forcing can be simply applied to its layer thickness after we have calculated the values of $W_0$ and $W_1$.

\paragraph{The one-layer model with moisture} This model has a one-layer setup. The physical quantities include:
\begin{itemize}
    \item The layer height $H$ and the column-averaged horizontal velocity $\mathbf{v}$.
    \item The column integrated water content in layer 1, $Q_1$, and the water content in the abyssal layer, $\alpha_0$.
    \item The local upward velocities at the bottom and top boundaries of this layer, $W_0$ and $W_1$.
\end{itemize}
Note that the above-listed variables imply several other quantities used (potentially as parameters) but not set as variables:
\begin{itemize}
    \item The potential temperature values $\theta_0$, $\theta_1$, and $\theta_2$, and gravity $g$.
    \item The time constants $\tau_1$ and $\tau_2$, controlling the precipitation and evaporation, respectively.
    \item The water contents in the abyssal layer $q_0$ and the saturated vapor concentration in the weather layer $q_s$.
\end{itemize}
And the equations that this shallow water model follows:
\begin{equation}
    \partial_t H + \nabla\cdot(H\mathbf{v})=W_0-W_1
\end{equation}
\begin{equation}
    \partial_t (H\mathbf{v})+(\mathbf{v}\cdot\nabla)(H\mathbf{v})+\mathbf{f}\times(H\mathbf{v})=g\frac{\theta_1}{\theta_0}\nabla z_0
\end{equation}
\begin{equation}
    W_1=\frac{LH(q-q_s)}{\tau_1 C_p(\theta_2-\theta_1)}
\end{equation}
\begin{equation}
    W_0=\frac{\alpha_0}{\tau_2 q_0}
\end{equation}
\begin{equation}
    \partial_t (Hq)+\nabla\cdot(Hq\mathbf{V}_i)=W_0 q_0-\frac{H(q-q_s)}{\tau_1}
\end{equation}
\begin{equation}
    \frac{d\alpha_0}{dt}=\frac{\int_A H(q-q_s)dA}{\tau_1\int_A dA} - \tau_2^{-1}\alpha_0
\end{equation}

\section{The two-layer coupled moist convection structures}
Thinking about Jupiter, we may want to go one step further to study how the ammonia distribution is coupled with the water convection dynamics.
\end{document}

